% ---
% id: reductions/rem:gateweight
% kind: rem
% title: "Weight of the licence in the consumer's gate"
% summary: >-
%   In untrusted-source mode, or where Alice's system has dimension one, the
%   consumer's mechanical check of the licensed class degenerates and its gate
%   rests entirely on the supplier's licence. The licence is therefore a
%   theorem with a named hypothesis rather than an attestation.
% prominence: standard
% status: accepted
% uses: [thm:licence]
% mentions: [meat-core/rem:cgate, req:licence-record]
% assumes: []
% refs: []
% provenance: newly-derived
% lean:
% ---

\begin{remark}
\label{rem:gateweight}
The consuming contract gates a key-length claim on a non-empty $\CList$ behind
two checks: (i) that the list lies in the licensed class of the declared source
mode, which is mechanical; and (ii) that the supplier's attack-set licence is
present, which no program can verify. It also records that check (i)
\emph{degenerates}: at $d_A=1$ every $1\times1$ Hermitian right-hand side is a
multiple of $\tau_A$, and in untrusted-source mode the type of an entry is
scalar by construction, so part (i) rejects nothing in untrusted-source mode
and nothing a producer cannot trivially avoid at $d_A=1$, and in both cases
the gate rests entirely on part (ii) (\xref{meat-core/rem:cgate}). This is the reason
Theorem~\ref{thm:licence} is stated as a theorem with an explicit hypothesis
rather than as an attestation: in the cases that matter, the attestation
\emph{is} the security argument, and the only defensible form for it is a
proof plus a named hypothesis that some other accepted contract discharges.
The \tref{licence record} names both, and the artefact that discharges the
hypothesis (Requirement~\ref{req:licence-record}).
\end{remark}
